% oct13-lem-contact: oct13.tex:644-649 (label lem:contact)
\paragraph{Paper excerpt.}
\begin{verbatim}
\begin{lemma}\label{lem:contact}
There exists $C(d,\drift,\gauge,p)<\infty$ such that, on the event of Section~\ref{sec:norm-event}, for large~$n$,
\begin{equation*}
H\leq Cr_nq_n=C\begin{cases}\sqrt{r_n\log n},&d=2,\\\log n,&d\geq3\end{cases}\, .
\end{equation*}
\end{lemma}
\end{verbatim}
\paragraph{Lean statement.}
\begin{verbatim}
theorem CERW.LimitShapes.contact_potential {d : ℕ} (hd : 2 ≤ d) :
    ∀ Ψ : EuclideanSpace ℝ (Fin d) → ℝ, CERW.IsNorm Ψ →
    ∀ ε : ℝ, 0 < ε → (∀ i : Fin d, ε * Ψ (CERW.coordVec i) < 1 / (d : ℝ)) →
    ∀ K : CERW.LimitShapes.EventConstants, 0 ≤ K.CCo → 0 ≤ K.CLoc → 0 ≤ K.CMart →
    ∃ C : ℝ, 0 < C ∧ ∃ n₀ : ℕ, 2 ≤ n₀ ∧
      ∀ ξ : Site d → EuclideanSpace ℝ (Fin d),
        (∀ x : Site d, x ≠ 0 → CERW.IsSubgradient Ψ (CERW.toSpace x) (ξ x)) → ξ 0 = 0 →
      ∀ (X : ℕ → Site d) (n : ℕ), n₀ ≤ n →
        CERW.Support.Norm.ContactEvent.PathTyping d X →
        CERW.LimitShapes.EventEstimates Ψ ε ξ K X n →
        ∀ y₀ : EuclideanSpace ℝ (Fin d), Ψ y₀ = CERW.normInnerRadius Ψ X n →
          y₀ ∈ closure (CERW.cellSet X n)ᶜ →
          CERW.normPotential d ε Ψ (CERW.cellSet X n) y₀ ≤
            C * CERW.Support.Norm.ContactEvent.scale d Ψ ε n *
              CERW.Support.Norm.ContactEvent.contactRate d Ψ ε n
\end{verbatim}
\paragraph{Proof of the registered statement.}
Fix the norm, drift and nonnegative event constants in the registered type.
The kernel gradient estimate bounds the actual predictable bracket of the
translated kernel by $C_g^2$ times its occupation-weighted sum; subtracting
the square of the conditional mean can only decrease this bracket.
Consequently the local actual-bracket constituent implies the weighted local
martingale estimate, with constant $K.\mathrm{CMart}(C_g+1)$.
The coarse constituent and the reduced local approximation are the other
inputs of the proved deterministic convolution estimate in
\texttt{ContactApproximation.lean}. Apply that estimate with these constants.
It yields a positive contact constant and a time threshold independent of
the subgradient field, path and contact point. Enlarge the threshold to at
least two and take the indicated three constituents of the sixfold event.
This proves the registered deterministic implication at every contact point.

This proof does not produce the event constants as functions of
$d,\kappa,\Psi,p$. The source's common-event probability assertion remains open,
so this node has conditional source coverage.
